% GENERATED FILE. Edit canonical lessons, then run npm run generate:books.
% canonical-source-hash: fnv1a64:1714a08c07a4e15f
% canonical-lessons: GU-C113-khissu, GU-C113-gol, GU-C113-madh, GU-C113-cana, GU-C113-mag

\chapter{Pocket, Jaggery, Honey, Chickpeas, Mung beans}
\label{ch:gu-pocket-jaggery-honey-chickpeas-mung-beans}
\hlchaptermodality{113}

\begin{chapteropening}
\textbf{By the end of this chapter:} \emph{I can say a pocket, jaggery, honey, chickpeas and mung beans.}

Complete the last lesson of chapter 113: i can say a pocket, jaggery, honey, chickpeas and mung beans.
\end{chapteropening}

\section[khissũ]{\gu{ખિસ્સું} (khissũ) — a pocket}

\label{lesson:GU-C113-khissu}



\begin{quote}
Before the new one: say the Gujarati for a sweater, then the Gujarati for a coat.
\end{quote}

\subsection*{You'll want to know: \gu{ખિસ્સું}}
\textbf{\gu{ખિસ્સું}} (\emph{khissũ}) — \textquotedblleft{}a pocket\textquotedblright{}.

\begin{grammarlens}[title={What you've built}]
One more word for this chapter.
\end{grammarlens}

\subsection*{Your turn}
\begin{itemize}
\raggedright
  \item \emph{Say it:} \emph{khissũ}
  \item \emph{Say it:} \emph{khissũ}, once more
  \item \emph{Say it:} say \emph{khissũ}
  \item \emph{Recall:} say the Gujarati for a sweater, then the Gujarati for a coat, then say \emph{khissũ} again
\end{itemize}

\begin{tcolorbox}[breakable,colback=teal!4,colframe=teal!35!black,title={Before you move on}]
What does \gu{ખિસ્સું} mean? (\textquotedblleft{}A pocket\textquotedblright{}.) Say it once more.
\end{tcolorbox}
\section[goḷ]{\gu{ગોળ} (goḷ) — jaggery}

\label{lesson:GU-C113-gol}



\begin{quote}
Before the new one: say the Gujarati for a coat, then the Gujarati for a pocket.
\end{quote}

\subsection*{You'll want to know: \gu{ગોળ}}
\textbf{\gu{ગોળ}} (\emph{goḷ}) — \textquotedblleft{}jaggery\textquotedblright{}.

\begin{grammarlens}[title={What you've built}]
One more word for this chapter.
\end{grammarlens}

\subsection*{Your turn}
\begin{itemize}
\raggedright
  \item \emph{Say it:} \emph{goḷ}
  \item \emph{Say it:} \emph{goḷ}, once more
  \item \emph{Say it:} say \emph{goḷ}
  \item \emph{Recall:} say the Gujarati for a coat, then the Gujarati for a pocket, then say \emph{goḷ} again
\end{itemize}

\begin{tcolorbox}[breakable,colback=teal!4,colframe=teal!35!black,title={Before you move on}]
What does \gu{ગોળ} mean? (\textquotedblleft{}Jaggery\textquotedblright{}.) Say it once more.
\end{tcolorbox}
\section[madh]{\gu{મધ} (madh) — honey}

\label{lesson:GU-C113-madh}



\begin{quote}
Before the new one: say the Gujarati for a pocket, then the Gujarati for jaggery.
\end{quote}

\subsection*{You'll want to know: \gu{મધ}}
\textbf{\gu{મધ}} (\emph{madh}) — \textquotedblleft{}honey\textquotedblright{}.

\begin{grammarlens}[title={What you've built}]
One more word for this chapter.
\end{grammarlens}

\subsection*{Your turn}
\begin{itemize}
\raggedright
  \item \emph{Say it:} \emph{madh}
  \item \emph{Say it:} \emph{madh}, once more
  \item \emph{Say it:} say \emph{madh}
  \item \emph{Recall:} say the Gujarati for a pocket, then the Gujarati for jaggery, then say \emph{madh} again
\end{itemize}

\begin{tcolorbox}[breakable,colback=teal!4,colframe=teal!35!black,title={Before you move on}]
What does \gu{મધ} mean? (\textquotedblleft{}Honey\textquotedblright{}.) Say it once more.
\end{tcolorbox}
\section[caṇā]{\gu{ચણા} (caṇā) — chickpeas}

\label{lesson:GU-C113-cana}



\begin{quote}
Before the new one: say the Gujarati for jaggery, then the Gujarati for honey.
\end{quote}

\subsection*{You'll want to know: \gu{ચણા}}
\textbf{\gu{ચણા}} (\emph{caṇā}) — \textquotedblleft{}chickpeas\textquotedblright{}.

\begin{grammarlens}[title={What you've built}]
One more word for this chapter.
\end{grammarlens}

\subsection*{Your turn}
\begin{itemize}
\raggedright
  \item \emph{Say it:} \emph{caṇā}
  \item \emph{Say it:} \emph{caṇā}, once more
  \item \emph{Say it:} say \emph{caṇā}
  \item \emph{Recall:} say the Gujarati for jaggery, then the Gujarati for honey, then say \emph{caṇā} again
\end{itemize}

\begin{tcolorbox}[breakable,colback=teal!4,colframe=teal!35!black,title={Before you move on}]
What does \gu{ચણા} mean? (\textquotedblleft{}Chickpeas\textquotedblright{}.) Say it once more.
\end{tcolorbox}
\section[mag]{\gu{મગ} (mag) — mung beans}

\label{lesson:GU-C113-mag}



\begin{quote}
Before the new one: say the Gujarati for honey, then the Gujarati for chickpeas.
\end{quote}

\subsection*{You'll want to know: \gu{મગ}}
\textbf{\gu{મગ}} (\emph{mag}) — \textquotedblleft{}mung beans\textquotedblright{}.

\begin{grammarlens}[title={What you've built}]
One more word for this chapter.
\end{grammarlens}

\subsection*{Your turn}
\begin{itemize}
\raggedright
  \item \emph{Say it:} \emph{mag}
  \item \emph{Say it:} \emph{mag}, once more
  \item \emph{Say it:} say \emph{mag}
  \item \emph{Recall:} say the Gujarati for honey, then the Gujarati for chickpeas, then say \emph{mag} again
\end{itemize}

\begin{tcolorbox}[breakable,colback=teal!4,colframe=teal!35!black,title={Before you move on}]
What does \gu{મગ} mean? (\textquotedblleft{}Mung beans\textquotedblright{}.) Say it once more.
\end{tcolorbox}
